\SourceLeaf{354}%
der entsprechenden unendlich kleinen Linien zu einander stehen, eine nach Verschiedenheit der Stellen in den Flächen veränderliche Zahl. In \textsc{Mercators} Projection z.\,B. ist die Vergrösserungszahl desto grösser, je entfernter vom Äquator, in der stereographischen Projection, je entfernter vom Augenpunkte die betreffenden Stellen sind.

Von jeder gegebenen Fläche sind auf einer andern gegebenen Fläche unendlich viele conforme Darstellungen möglich; die allgemeine Auflösung umfasst sie sämmtlich, indem sie eine arbiträre Function enthält, welche nach Gefallen oder den jedesmaligen Zwecken gemäss bestimmt werden kann. Wenn nur ein Theil der einen Fläche übertragen werden soll, ist es in der Regel am vortheilhaftesten, eine solche conforme Darstellung zu wählen, bei welcher innerhalb der darzustellenden Fläche die Ungleichheiten des Vergrösserungsverhältnisses in den möglich engsten Grenzen bleiben.

Die Aufgabe der conformen Übertragung der Ellipsoidfläche auf die Kugelfläche ist in der angeführten Schrift unter den Beispielen besonders abgehandelt, und der allgemeinen Auflösung sind zwei specielle beigefügt, wovon die eine vorzugsweise für die Darstellung der ganzen Ellipsoidfläche geeignet, die andere hingegen weit zweckmässiger ist, wenn (wie es immer bei bestimmten Anwendungen auf die Geodaesie der Fall ist) nur ein mässiger Theil der als ellipsoidisch betrachteten Erdfläche auf eine Kugelfläche conform übertragen werden soll. In wie hohem Grade diese zweite Darstellungsart der oben ausgesprochenen Forderung genügt, ist aus einem a.\,a.\,O. aufgestellten Beispiele abzunehmen, wo die Veränderlichkeit des Vergrösserungsverhältnisses innerhalb einer Zone von fünf Breitengraden nur \(\frac1{530000}\) beträgt. Es sind ferner daselbst die Hauptzüge der Methode, wie überhaupt eine conforme Übertragung zur Berechnung eines Dreieckssystems benutzt werden kann, im Allgemeinen angedeutet, die eigentliche Ausführung aber, und die Anwendung auf diese bestimmte Übertragungsart einer späteren Bearbeitung vorbehalten.

Die gegenwärtige Abhandlung ist nun dazu bestimmt, diese Verpflichtung auszulösen, obwohl nicht ganz in derselben Art, wie sie eingegangen war: es wird nemlich darin nicht die eben erwähnte, sondern eine davon verschiedene dritte specielle Auflösung der Aufgabe zum Grunde gelegt, durch welche der beabsichtigte Zweck noch vollkommener erreicht wird. In diesen Blättern müssen wir uns damit begnügen, nur im Allgemeinen einen Begriff davon zu geben.
